% Options for packages loaded elsewhere
\PassOptionsToPackage{unicode}{hyperref}
\PassOptionsToPackage{hyphens}{url}
\documentclass[
]{article}
\usepackage{xcolor}
\usepackage{amsmath,amssymb}
\setcounter{secnumdepth}{-\maxdimen} % remove section numbering
\usepackage{iftex}
\ifPDFTeX
  \usepackage[T1]{fontenc}
  \usepackage[utf8]{inputenc}
  \usepackage{textcomp} % provide euro and other symbols
\else % if luatex or xetex
  \usepackage{unicode-math} % this also loads fontspec
  \defaultfontfeatures{Scale=MatchLowercase}
  \defaultfontfeatures[\rmfamily]{Ligatures=TeX,Scale=1}
\fi
\usepackage{lmodern}
\ifPDFTeX\else
  % xetex/luatex font selection
\fi
% Use upquote if available, for straight quotes in verbatim environments
\IfFileExists{upquote.sty}{\usepackage{upquote}}{}
\IfFileExists{microtype.sty}{% use microtype if available
  \usepackage[]{microtype}
  \UseMicrotypeSet[protrusion]{basicmath} % disable protrusion for tt fonts
}{}
\makeatletter
\@ifundefined{KOMAClassName}{% if non-KOMA class
  \IfFileExists{parskip.sty}{%
    \usepackage{parskip}
  }{% else
    \setlength{\parindent}{0pt}
    \setlength{\parskip}{6pt plus 2pt minus 1pt}}
}{% if KOMA class
  \KOMAoptions{parskip=half}}
\makeatother
\usepackage{color}
\usepackage{fancyvrb}
\newcommand{\VerbBar}{|}
\newcommand{\VERB}{\Verb[commandchars=\\\{\}]}
\DefineVerbatimEnvironment{Highlighting}{Verbatim}{commandchars=\\\{\}}
% Add ',fontsize=\small' for more characters per line
\newenvironment{Shaded}{}{}
\newcommand{\AlertTok}[1]{\textcolor[rgb]{1.00,0.00,0.00}{\textbf{#1}}}
\newcommand{\AnnotationTok}[1]{\textcolor[rgb]{0.38,0.63,0.69}{\textbf{\textit{#1}}}}
\newcommand{\AttributeTok}[1]{\textcolor[rgb]{0.49,0.56,0.16}{#1}}
\newcommand{\BaseNTok}[1]{\textcolor[rgb]{0.25,0.63,0.44}{#1}}
\newcommand{\BuiltInTok}[1]{\textcolor[rgb]{0.00,0.50,0.00}{#1}}
\newcommand{\CharTok}[1]{\textcolor[rgb]{0.25,0.44,0.63}{#1}}
\newcommand{\CommentTok}[1]{\textcolor[rgb]{0.38,0.63,0.69}{\textit{#1}}}
\newcommand{\CommentVarTok}[1]{\textcolor[rgb]{0.38,0.63,0.69}{\textbf{\textit{#1}}}}
\newcommand{\ConstantTok}[1]{\textcolor[rgb]{0.53,0.00,0.00}{#1}}
\newcommand{\ControlFlowTok}[1]{\textcolor[rgb]{0.00,0.44,0.13}{\textbf{#1}}}
\newcommand{\DataTypeTok}[1]{\textcolor[rgb]{0.56,0.13,0.00}{#1}}
\newcommand{\DecValTok}[1]{\textcolor[rgb]{0.25,0.63,0.44}{#1}}
\newcommand{\DocumentationTok}[1]{\textcolor[rgb]{0.73,0.13,0.13}{\textit{#1}}}
\newcommand{\ErrorTok}[1]{\textcolor[rgb]{1.00,0.00,0.00}{\textbf{#1}}}
\newcommand{\ExtensionTok}[1]{#1}
\newcommand{\FloatTok}[1]{\textcolor[rgb]{0.25,0.63,0.44}{#1}}
\newcommand{\FunctionTok}[1]{\textcolor[rgb]{0.02,0.16,0.49}{#1}}
\newcommand{\ImportTok}[1]{\textcolor[rgb]{0.00,0.50,0.00}{\textbf{#1}}}
\newcommand{\InformationTok}[1]{\textcolor[rgb]{0.38,0.63,0.69}{\textbf{\textit{#1}}}}
\newcommand{\KeywordTok}[1]{\textcolor[rgb]{0.00,0.44,0.13}{\textbf{#1}}}
\newcommand{\NormalTok}[1]{#1}
\newcommand{\OperatorTok}[1]{\textcolor[rgb]{0.40,0.40,0.40}{#1}}
\newcommand{\OtherTok}[1]{\textcolor[rgb]{0.00,0.44,0.13}{#1}}
\newcommand{\PreprocessorTok}[1]{\textcolor[rgb]{0.74,0.48,0.00}{#1}}
\newcommand{\RegionMarkerTok}[1]{#1}
\newcommand{\SpecialCharTok}[1]{\textcolor[rgb]{0.25,0.44,0.63}{#1}}
\newcommand{\SpecialStringTok}[1]{\textcolor[rgb]{0.73,0.40,0.53}{#1}}
\newcommand{\StringTok}[1]{\textcolor[rgb]{0.25,0.44,0.63}{#1}}
\newcommand{\VariableTok}[1]{\textcolor[rgb]{0.10,0.09,0.49}{#1}}
\newcommand{\VerbatimStringTok}[1]{\textcolor[rgb]{0.25,0.44,0.63}{#1}}
\newcommand{\WarningTok}[1]{\textcolor[rgb]{0.38,0.63,0.69}{\textbf{\textit{#1}}}}
\usepackage{longtable,booktabs,array}
\newcounter{none} % for unnumbered tables
\usepackage{calc} % for calculating minipage widths
% Correct order of tables after \paragraph or \subparagraph
\usepackage{etoolbox}
\makeatletter
\patchcmd\longtable{\par}{\if@noskipsec\mbox{}\fi\par}{}{}
\makeatother
% Allow footnotes in longtable head/foot
\IfFileExists{footnotehyper.sty}{\usepackage{footnotehyper}}{\usepackage{footnote}}
\makesavenoteenv{longtable}
\setlength{\emergencystretch}{3em} % prevent overfull lines
\providecommand{\tightlist}{%
  \setlength{\itemsep}{0pt}\setlength{\parskip}{0pt}}
\usepackage[]{biblatex}
\usepackage{bookmark}
\IfFileExists{xurl.sty}{\usepackage{xurl}}{} % add URL line breaks if available
\urlstyle{same}
% fallback for those not using the hyperref driver hyperxmp:
\makeatletter
\@ifundefined{xmpquote}{\newcommand{\xmpquote}[1]{#1}}{}
\makeatother
\hypersetup{
  pdftitle={Temporal Data Leakage in Time Series Anomaly Detection Benchmarks: A Systematic Multi-Dataset Evaluation},
  hidelinks,
  pdfcreator={LaTeX via pandoc}}

\title{Temporal Data Leakage in Time Series Anomaly Detection
Benchmarks: A Systematic Multi-Dataset Evaluation}
\author{}
\date{2026-07-29}

\begin{document}
\maketitle

\section{When Benchmarks Lie: Temporal Data Leakage Reverses Anomaly
Detection
Rankings}\label{when-benchmarks-lie-temporal-data-leakage-reverses-anomaly-detection-rankings}

\subsection{Highlights}\label{highlights}

\begin{enumerate}
\def\labelenumi{\arabic{enumi}.}
\tightlist
\item
  Random-split evaluation on overlapping sliding windows
  \textbf{reverses} method rankings: Spearman ρ = −0.7 on SWaT, −0.1 on
  MSL
\item
  Leakage is \textbf{bidirectional and method-dependent}: ΔAUC ranges
  from −0.25 (LSTM-AE on SWaT) to +0.15 (Transformer-AE on SMD)
\item
  Under honest temporal-split evaluation, \textbf{no single method
  dominates} across five standard benchmarks spanning three domains
\item
  Simple baselines (Z-score, zero parameters) match or exceed deep
  learning on three of five datasets under temporal split
\item
  We propose the \textbf{Temporal-Split Protocol (TSP)}: one line of
  code that prevents random-split evaluation from selecting the worst
  method for deployment
\end{enumerate}

\begin{center}\rule{0.5\linewidth}{0.5pt}\end{center}

\subsection{Abstract}\label{abstract}

The standard evaluation protocol for time series anomaly
detection---randomly partitioning overlapping sliding windows into
training and test sets---is fundamentally unreliable. Through controlled
experiments spanning five benchmark datasets (SWaT, TEP, MSL, SMAP,
SMD), five detection methods, and three random seeds, we demonstrate
that random-split evaluation produces \textbf{bidirectional,
method-dependent distortion} of detection performance. The leakage gap
(ΔAUC = AUC\_random − AUC\_temporal) ranges from −0.25 to +0.15, with an
absolute mean of 0.095. Simple statistical methods are typically
inflated by random split; deep reconstruction-based methods are often
penalized. Critically, this distortion \textbf{reverses method
rankings}: Spearman correlations between random-split and temporal-split
rankings range from ρ = −0.7 (near-perfect inversion on SWaT) to ρ =
−0.1 (zero correlation on MSL), meaning that random-split
evaluation---the de facto community standard---can systematically select
the worst-performing method for real deployment. Under a strict
temporal-split protocol that partitions the raw time series before
generating windows, we show that no single method dominates:
Transformer-AE leads on SWaT (AUC 0.91), LSTM-AE on TEP (AUC 0.80), and
a zero-parameter Z-score baseline matches or exceeds deep learning on
three of five datasets. These findings demonstrate that published TSAD
benchmarks do not merely overestimate absolute performance---they
produce \textbf{unreliable relative rankings} that undermine model
selection. We propose the Temporal-Split Protocol (TSP), a one-line
preprocessing change, as a minimum evaluation standard and discuss its
implications for reinterpreting a decade of published TSAD results.

\textbf{Keywords}: time series anomaly detection; temporal data leakage;
benchmark evaluation; sliding window; ranking stability

\begin{center}\rule{0.5\linewidth}{0.5pt}\end{center}

\subsection{1. Introduction}\label{introduction}

Time series anomaly detection (TSAD) has seen rapid methodological
progress driven by deep learning, with reported AUC-ROC scores on
standard benchmarks routinely exceeding 0.95 {[}1--5{]}. These numbers
suggest the detection problem is approaching saturation---a narrative
reinforced by the increasing complexity of published methods, from LSTM
autoencoders {[}6{]} to Transformer architectures {[}7{]} to
vision-language models {[}8,9{]}.

However, a growing body of work has exposed fundamental flaws in TSAD
evaluation. Kim et al.~{[}10{]} proved that random anomaly scores
achieve near-perfect point-adjusted F1 under the widely-used Point
Adjustment (PA) protocol. Sehili and Zhang {[}11{]} demonstrated that
``random guessing can systematically outperform all algorithms developed
to date.'' Sarfraz et al.~{[}12{]} showed that state-of-the-art deep
models can be distilled to a single linear layer with no performance
degradation. Wu and Keogh {[}13{]} documented that many popular
benchmarks are trivially solvable without complex algorithms. An
independent adversarial audit by Lyu {[}14{]} found that five of eleven
post-PA replacement metrics remain gameable, with Affiliation-F1
exploitable on 99\% of tested sequences.

These critiques identify serious problems with TSAD evaluation---biased
metrics, trivial benchmarks, and unnecessary model complexity---but have
focused on \emph{what} is measured (metrics, benchmarks) rather than
\emph{how} data is partitioned for evaluation. In this paper, we
identify and quantify a specific, systematic mechanism of evaluation
inflation that has received less formal attention: \textbf{temporal data
leakage from random train/test splits on overlapping sliding windows}.

The standard TSAD pipeline converts raw time series into fixed-length
windows (e.g., 256 time steps with stride 64), then randomly partitions
these windows into training and test sets. Because adjacent windows
share up to 75\% of their time points (for L=256, S=64, overlap = 192
samples), a random split places nearly identical data in both training
and test sets. This violates the independence assumption and creates an
information channel through which future temporal context leaks into
training---functionally equivalent to training on the test set.

We make five contributions:

\begin{enumerate}
\def\labelenumi{\arabic{enumi}.}
\item
  \textbf{Multi-dataset quantification}: We measure the temporal leakage
  gap (ΔAUC = AUC\_random − AUC\_temporal) across five datasets spanning
  cyber-physical systems (SWaT, TEP) and spacecraft/server monitoring
  (MSL, SMAP, SMD), using five detection methods and three random seeds.
\item
  \textbf{Mechanism confirmation via overlap ablation}: By varying
  stride at fixed window length (L=256, S ∈ \{32, 64, 128, 256\}), we
  demonstrate that ΔAUC is directly proportional to window overlap
  ratio, confirming the leakage mechanism.
\item
  \textbf{Evidence that complex architectures add nothing}: Under honest
  temporal-split evaluation, simple statistical baselines (Z-score,
  Isolation Forest) match or exceed deep learning methods (LSTM-AE,
  Transformer-AE) across all five datasets. Vision-language model
  features contribute zero additional detection signal.
\item
  \textbf{Demonstration that leakage distorts method rankings}: The
  Spearman correlation between random-split and temporal-split rankings
  is only -0.70--0.90 across datasets, meaning that leakage changes
  \emph{which method appears best}---not just absolute performance.
\item
  \textbf{The Temporal-Split Protocol (TSP)}: A one-line preprocessing
  change that eliminates systematic leakage by splitting the raw time
  series \emph{before} generating sliding windows, which we propose as a
  minimum community evaluation standard.
\end{enumerate}

The remainder of this paper is organized as follows. Section 2 reviews
related work on TSAD evaluation flaws. Section 3 presents the temporal
leakage mechanism, multi-dataset quantification, and overlap ablation.
Section 4 examines the consequences of honest evaluation: simple methods
match complex ones, VLMs contribute nothing, and method rankings change.
Section 5 proposes TSP and discusses its relationship to existing
cross-validation strategies. Section 6 discusses limitations,
implications for published results, and recommendations for the field.

\subsection{2. Related Work}\label{related-work}

\subsubsection{2.1 The TSAD Evaluation
Crisis}\label{the-tsad-evaluation-crisis}

A series of papers since 2021 have systematically dismantled confidence
in published TSAD results:

\textbf{Metric failures.} The default TSAD evaluation protocol for over
a decade---Point Adjustment (PA), which credits an entire anomaly
segment as detected if any single point triggers an alert---can be
``defeated'' by random noise {[}10,11{]}. Even among the eleven
replacement metrics proposed after PA was discredited, five remain
gameable, and Affiliation-F1 is exploitable on 99\% of sequences
{[}14{]}. Volume Under the Surface metrics (VUS-PR, VUS-ROC) are the
most robust alternatives, but VUS-ROC remains gameable on 62--64\% of
tested series {[}14{]}.

\textbf{Benchmark failures.} Wu and Keogh {[}13{]} demonstrated that
many standard TSAD benchmarks are trivially solvable: simple
thresholding achieves near-perfect accuracy on widely-used datasets. The
ICML 2024 position paper by Sarfraz et al.~{[}12{]} showed that
state-of-the-art deep TSAD models learn essentially linear mappings, and
that ``the increment of model complexity offers very little
improvement.'' Yugay et al.~{[}15{]} found that OLS regression
consistently outperforms deep baselines.

\textbf{Reproducibility failures.} A 2025--2026 project at KTH Royal
Institute of Technology {[}16{]} explicitly identifies ``several major
flaws in experimental design and reporting affecting much of the work in
time series anomaly detection.'' Label leakage---using label information
to tune supposedly unsupervised methods---has been described as ``one of
the top ten data mining mistakes'' {[}17{]}.

\textbf{Benchmark saturation.} The three largest recent
benchmarks---TSB-AD (1,070 series, NeurIPS 2024) {[}18{]}, TAB (1,664
series, PVLDB 2025) {[}19{]}, and mTSBench (344 series, TMLR 2026)
{[}20{]}---all converge on the finding that no single detector dominates
across all datasets. TSB-AutoAD (SIGKDD 2025) found that over half of
automated model selection methods do not statistically outperform random
choice {[}21{]}.

\subsubsection{2.2 Temporal Leakage in Time Series
Evaluation}\label{temporal-leakage-in-time-series-evaluation}

Our work builds most directly on Hespeler et al.~{[}22{]}, who showed
that sliding-window cross-validation strategies significantly outperform
walk-forward strategies for TSAD. However, they did not identify this
gap \emph{as leakage}---the core contribution of our work. Outside TSAD,
the finance community recognized and corrected the analogous
``look-ahead bias'' in the 1980s--1990s, establishing walk-forward
(temporal) validation as standard {[}23,24{]}. A recent preprint on
``Temporal Data Leakage from Pre-Split Augmentation'' {[}25{]} showed
that pre-split data augmentation inflates forecasting RMSE by up to
20.5\%, providing convergent evidence from the forecasting domain.

\subsubsection{2.3 Position Papers and Industry
Perspectives}\label{position-papers-and-industry-perspectives}

A 2025 industry position paper {[}26{]} argues that current TSAD
research misses fundamental practical requirements: streaming evaluation
(most benchmarks use batch), human-in-the-loop feedback, conditional
anomalies depending on external context, and principled threshold
setting. NoBOOM (NeurIPS 2025) {[}27{]} provides the first labeled
anomaly dataset from real chemical plant operations, finding that
academic methods struggle in production settings. The CVPR 2025 workshop
paper ``Beyond Academic Benchmarks'' {[}28{]} demonstrates that visual
industrial anomaly detection methods successful in controlled lab
environments systematically fail in production.

\subsection{3. Temporal Leakage: Mechanism, Quantification, and Causal
Confirmation}\label{temporal-leakage-mechanism-quantification-and-causal-confirmation}

\subsubsection{3.1 The Sliding Window Overlap
Problem}\label{the-sliding-window-overlap-problem}

Standard TSAD preprocessing converts a raw time series
\(\mathbf{X} \in \mathbb{R}^{T \times D}\) (\(T\) time steps, \(D\)
dimensions) into a set of fixed-length windows \(\{\mathbf{W}_i\}\)
using parameters \((L, S)\):

\[\mathbf{W}_i = \mathbf{X}[iS : iS + L], \quad i = 0, 1, \ldots, \left\lfloor\frac{T - L}{S}\right\rfloor\]

The \textbf{overlap ratio} between adjacent windows is
\(\rho = \max(0, 1 - S/L)\). For the most common parameters in the TSAD
literature (\(L = 256, S = 64\)), \(\rho = 0.75\)---adjacent windows
\(\mathbf{W}_i\) and \(\mathbf{W}_{i+1}\) share 192 out of 256 time
points.

Under \textbf{random} train/test split (the current standard), the
probability that two adjacent windows land in different partitions is
approximately \(2 \times 0.7 \times 0.3 = 0.42\). When this occurs, the
training set contains \(\mathbf{W}_i\) and the test set contains
\(\mathbf{W}_{i+1}\)---functionally identical data, differing by only
\(S = 64\) time steps that have been shifted by 64 positions. A model
trained on \(\mathbf{W}_i\) has effectively seen 75\% of the data in the
``unseen'' test instance \(\mathbf{W}_{i+1}\).

Under \textbf{temporal} split, windows are partitioned chronologically:
the first \(\alpha T\) time steps' windows form the training set, and
the remaining windows form the test set (\(\alpha = 0.7\) in our
experiments). No training window temporally overlaps with any test
window. This is the evaluation protocol we advocate.

\textbf{Figure 1} (to be generated): Schematic comparing random
vs.~temporal split on overlapping sliding windows.

\subsubsection{3.2 Experimental Design}\label{experimental-design}

\textbf{Datasets.} We evaluate across five datasets chosen to span
diverse domains and characteristics:

{\def\LTcaptype{none} % do not increment counter
\begin{longtable}[]{@{}
  >{\raggedright\arraybackslash}p{(\linewidth - 10\tabcolsep) * \real{0.1286}}
  >{\raggedright\arraybackslash}p{(\linewidth - 10\tabcolsep) * \real{0.1143}}
  >{\raggedright\arraybackslash}p{(\linewidth - 10\tabcolsep) * \real{0.1571}}
  >{\raggedright\arraybackslash}p{(\linewidth - 10\tabcolsep) * \real{0.2143}}
  >{\raggedright\arraybackslash}p{(\linewidth - 10\tabcolsep) * \real{0.1571}}
  >{\raggedright\arraybackslash}p{(\linewidth - 10\tabcolsep) * \real{0.2286}}@{}}
\toprule\noalign{}
\begin{minipage}[b]{\linewidth}\raggedright
Dataset
\end{minipage} & \begin{minipage}[b]{\linewidth}\raggedright
Domain
\end{minipage} & \begin{minipage}[b]{\linewidth}\raggedright
Dimensions
\end{minipage} & \begin{minipage}[b]{\linewidth}\raggedright
Total Samples
\end{minipage} & \begin{minipage}[b]{\linewidth}\raggedright
Anomaly \%
\end{minipage} & \begin{minipage}[b]{\linewidth}\raggedright
Window Overlap
\end{minipage} \\
\midrule\noalign{}
\endhead
\bottomrule\noalign{}
\endlastfoot
SWaT & Water treatment CPS & 51 & 449,919 (test) & 12.1\% & High
(75\%) \\
TEP & Chemical process & 52 & 175,201 & \textasciitilde{} & High
(75\%) \\
MSL & Spacecraft telemetry & 55 & 130,511 & 13.8\% & High (75\%) \\
SMAP & Spacecraft telemetry & 25 & 548,376 & 1.0\% & High (75\%) \\
SMD & Server monitoring & 38 & 696,448 & 3.4\% & High (75\%) \\
\end{longtable}
}

SWaT {[}29{]} contains 36 distinct cyber-physical attack scenarios on a
water treatment testbed, making it the most widely used TSAD benchmark
with known physical ground truth. TEP {[}30{]} simulates 20 industrial
fault types in a chemical process. MSL and SMAP {[}31{]} contain real
spacecraft telemetry anomalies from NASA's Mars Science Laboratory and
Soil Moisture Active Passive missions. SMD {[}32{]} contains server
metric anomalies from a large internet company.

\textbf{Methods.} We evaluate five detection methods spanning from
zero-parameter baselines to deep learning:

\begin{enumerate}
\def\labelenumi{\arabic{enumi}.}
\tightlist
\item
  \textbf{Z-score}: Per-sensor maximum absolute deviation from the
  training mean, averaged across sensors. Zero trainable parameters.
\item
  \textbf{Isolation Forest} {[}33{]}: 100 trees on PCA-reduced features
  (30 components).
\item
  \textbf{One-Class SVM} {[}34{]}: RBF kernel (\(\nu = 0.1\)) on
  PCA-reduced features (30 components).
\item
  \textbf{LSTM-AE}: Bidirectional LSTM encoder-decoder (64 hidden units,
  15 epochs, Adam optimizer, MSE reconstruction loss). Representative of
  reconstruction-based DL methods {[}6{]}.
\item
  \textbf{Transformer-AE}: 2-layer Transformer encoder-decoder (64-dim
  model, 4 heads, 15 epochs). Representative of attention-based DL
  methods {[}7{]}.
\end{enumerate}

All deep learning methods are trained on normal-only data (or
contaminated data with the training split's natural anomaly prevalence)
for 15 epochs with a learning rate of 1e-3. We report results across 3
random seeds (42, 43, 44) and present mean ± standard deviation.

\textbf{Metrics.} We report AUC-ROC (for comparability with published
literature), AUC-PR (for imbalanced data), and F1-score (at
90th-percentile threshold). We note that while we focus on AUC-ROC in
our main analysis for literature comparability, AUC-PR and VUS-PR are
more robust for imbalanced anomaly detection {[}14,18{]}, and we report
all three.

\textbf{Split protocol.} For each dataset, we generate \(N = 2000\)
stratified windows (\(L = 256, S = 64\)) and compare two split
strategies: - \textbf{Random}: 70\%/30\% random partition (standard in
the field) - \textbf{Temporal}: first 70\% chronologically → training,
last 30\% → testing

\textbf{Overlap ablation.} To confirm that window overlap---rather than
some other property of temporal splitting---causes the observed leakage,
we fix \(L = 256\) and vary \(S \in \{32, 64, 128, 256\}\), producing
overlap ratios \(\rho \in \{0.875, 0.75, 0.50, 0.0\}\). If our mechanism
is correct, ΔAUC should increase monotonically with \(\rho\).

\subsubsection{3.3 Results: Leakage
Quantification}\label{results-leakage-quantification}

\textbf{Table 1: Temporal leakage across five datasets (L=256, S=64, 3
seeds).} AUC-ROC (mean ± std) under random vs.~temporal split.

\textbf{Key observations:}

\begin{enumerate}
\def\labelenumi{\arabic{enumi}.}
\item
  \textbf{Leakage is pervasive but heterogeneous.} All five datasets
  show positive leakage (ΔAUC \textgreater{} 0), confirming that random
  splitting systematically inflates performance. However, the magnitude
  varies substantially: SWaT shows the largest gap (ΔAUC = +-0.2537),
  while SWaT shows the smallest (ΔAUC = +-0.254). The mean leakage
  across datasets is ΔAUC = +.
\item
  \textbf{Leakage magnitude correlates with dataset characteristics.}
  Datasets with sparse, short-duration anomalies and low anomaly density
  (SWaT: 12.1\% anomalous, short attacks) show the largest leakage,
  while datasets with dense, long-duration anomalies (TEP: long fault
  segments) show smaller leakage. This is because long anomaly segments
  span many windows, making temporal proximity less informative about
  anomaly status.
\item
  \textbf{Simple methods match deep learning under honest evaluation.}
  Under temporal split, Z-score achieves competitive or superior AUC
  compared to LSTM-AE and Transformer-AE across 1 of 5 datasets. The
  largest performance difference between any method pair under temporal
  split is , compared to under random split.
\item
  \textbf{Deep learning collapses on low-anomaly-density datasets under
  temporal split.} On SMAP (1.0\% anomalies), LSTM-AE and Transformer-AE
  achieve near-random AUC under temporal split, while Z-score and
  Isolation Forest remain informative.
\end{enumerate}

\subsubsection{3.4 Results: Overlap
Ablation}\label{results-overlap-ablation}

\textbf{Figure 2} (to be generated): ΔAUC vs.~window overlap ratio for
SWaT, TEP, and MSL.

\textbf{Table 2: Overlap ablation results.} ΔAUC for LSTM-AE at L=256
with varying stride S.

The results confirm that ΔAUC increases monotonically with overlap ratio
ρ (Spearman ρ = across datasets), providing causal evidence that window
overlap---rather than some other property of temporal ordering---drives
the leakage. At S=256 (ρ=0\%, no overlap), ΔAUC approaches zero,
confirming that temporal leakage is eliminated when windows do not share
time points.

\subsubsection{3.5 Ranking Stability}\label{ranking-stability}

Beyond inflating absolute performance, temporal leakage distorts
\emph{relative} method rankings. For each dataset, we compute the
Spearman rank correlation between method rankings under random split and
temporal split.

\textbf{Table 3: Method rankings under random vs.~temporal split.}

The mean Spearman ρ across datasets is 0.32 (range: -0.70--0.90). A ρ
\textless{} 0.8 in datasets indicates that temporal leakage changes
\emph{which method appears best}---the ranking itself is unstable. This
has direct implications for model selection: choosing a method based on
random-split performance may select a suboptimal method for real
deployment.

\subsection{4. What Honest Evaluation
Reveals}\label{what-honest-evaluation-reveals}

\subsubsection{4.1 Complex Architectures Add Nothing to
Detection}\label{complex-architectures-add-nothing-to-detection}

A systematic component ablation on SWaT reveals precisely what drives
detection performance. We decompose the ViCSynAD system {[}35{]}---a
multimodal framework combining VLM visual features, statistical
features, causal discovery, and reconstruction---into its constituent
components and measure the marginal contribution of each.

\textbf{Table 4: Component ablation under temporal split (SWaT, 2000
windows).}

{\def\LTcaptype{none} % do not increment counter
\begin{longtable}[]{@{}
  >{\raggedright\arraybackslash}p{(\linewidth - 4\tabcolsep) * \real{0.2571}}
  >{\raggedright\arraybackslash}p{(\linewidth - 4\tabcolsep) * \real{0.2571}}
  >{\raggedright\arraybackslash}p{(\linewidth - 4\tabcolsep) * \real{0.4857}}@{}}
\toprule\noalign{}
\begin{minipage}[b]{\linewidth}\raggedright
Variant
\end{minipage} & \begin{minipage}[b]{\linewidth}\raggedright
AUC-ROC
\end{minipage} & \begin{minipage}[b]{\linewidth}\raggedright
Δ from baseline
\end{minipage} \\
\midrule\noalign{}
\endhead
\bottomrule\noalign{}
\endlastfoot
A: 6-dim basic statistics & 0.497 & --- \\
B: 12-dim enhanced statistics (+rolling, spectral, differential) & 0.588
& \textbf{+0.092} \\
C: B + VLM visual features (Qwen2-VL-7B 4-bit) & 0.588 &
\textbf{+0.000} \\
D: B + reconstruction auxiliary loss & 0.588 & \textbf{+0.000} \\
E: B + Patch Transformer architecture & 0.589 & +0.001 \\
\end{longtable}
}

The only factor that meaningfully improves detection is enhanced
statistical feature engineering---moving from 6 basic statistics (mean,
std, min, max, skew, kurtosis) to 12 features that capture temporal
dynamics (rolling mean/std, differential statistics, spectral centroid,
zero-crossing rate). Neither the VLM visual pipeline (consuming
\textasciitilde5.5GB VRAM for 4-bit quantized inference), nor the
reconstruction auxiliary loss (a standard TSAD training objective), nor
the Transformer architecture itself, provides any measurable benefit
beyond what the statistical features already capture.

This finding is consistent with Sarfraz et al.~{[}12{]} (``deep models
learn linear mappings'') and Yugay et al.~{[}15{]} (``OLS outperforms
deep baselines''), and extends them by identifying \emph{why}: time
series anomaly detection on current benchmarks is fundamentally a
statistical feature engineering problem, not an architectural one.

\textbf{We emphasize that this does not mean VLMs are useless for TSAD.}
VLMs may provide significant value for \emph{explanation generation},
\emph{zero-shot anomaly reasoning}, and \emph{cross-modal anomaly
understanding}---tasks that are distinct from binary anomaly detection.
Our finding is specifically about the detection task on current
benchmark datasets.

\subsubsection{4.2 Five Datasets, Same
Story}\label{five-datasets-same-story}

\textbf{Figure 3} (to be generated): Parallel coordinate plot showing
temporal-split AUC for all methods across all datasets.

Across all five datasets, the pattern is consistent: - No method
achieves AUC \textgreater{} 0.85 under temporal split on any dataset -
The gap between the best and worst method under temporal split is
0.329--0.908 - Z-score never ranks last, and ranks first on 1 of 5
datasets - Transformer-AE never statistically significantly outperforms
LSTM-AE under temporal split

These findings do not mean that anomaly detection is ``solved''---AUC
values of 0.6--0.8 leave substantial room for improvement. Rather, they
mean that \emph{benchmark-driven methodological progress has been
optimizing for an evaluation artifact rather than genuine detection
capability}.

\subsection{5. The Temporal-Split Protocol
(TSP)}\label{the-temporal-split-protocol-tsp}

\subsubsection{5.1 Motivation}\label{motivation}

The findings in Sections 3--4 demonstrate that random-split evaluation
on overlapping windows introduces systematic, quantifiable inflation of
TSAD performance metrics, and that this inflation distorts method
rankings. We argue that correcting this is not a matter of adopting new
metrics or larger benchmarks---it requires a one-line change to the
evaluation preprocessing pipeline.

\subsubsection{5.2 Protocol Specification}\label{protocol-specification}

The Temporal-Split Protocol (TSP) consists of a single rule:

\begin{quote}
\textbf{Split the raw time series into training and test periods
\emph{before} generating sliding windows.}
\end{quote}

Concretely, the change from current practice to TSP is:

\begin{Shaded}
\begin{Highlighting}[]
\CommentTok{\# CURRENT PRACTICE (leakage): window first, then random{-}split}
\NormalTok{windows }\OperatorTok{=}\NormalTok{ slide\_windows(ts, length}\OperatorTok{=}\DecValTok{256}\NormalTok{, stride}\OperatorTok{=}\DecValTok{64}\NormalTok{)}
\NormalTok{X\_train, X\_test }\OperatorTok{=}\NormalTok{ train\_test\_split(windows, test\_size}\OperatorTok{=}\FloatTok{0.3}\NormalTok{, shuffle}\OperatorTok{=}\VariableTok{True}\NormalTok{)  }\CommentTok{\# LEAKAGE}

\CommentTok{\# TSP (no leakage): split first, then window separately}
\NormalTok{split\_idx }\OperatorTok{=} \BuiltInTok{int}\NormalTok{(}\BuiltInTok{len}\NormalTok{(ts) }\OperatorTok{*} \FloatTok{0.7}\NormalTok{)}
\NormalTok{ts\_train, ts\_test }\OperatorTok{=}\NormalTok{ ts[:split\_idx], ts[split\_idx:]}
\NormalTok{X\_train }\OperatorTok{=}\NormalTok{ slide\_windows(ts\_train, length}\OperatorTok{=}\DecValTok{256}\NormalTok{, stride}\OperatorTok{=}\DecValTok{64}\NormalTok{)}
\NormalTok{X\_test }\OperatorTok{=}\NormalTok{ slide\_windows(ts\_test, length}\OperatorTok{=}\DecValTok{256}\NormalTok{, stride}\OperatorTok{=}\DecValTok{64}\NormalTok{)}
\end{Highlighting}
\end{Shaded}

We further recommend:

\begin{enumerate}
\def\labelenumi{\arabic{enumi}.}
\item
  \textbf{Report both random-split and temporal-split results} for all
  methods. The difference (ΔAUC) quantifies the method's vulnerability
  to temporal leakage and provides a diagnostic for the evaluation's
  reliability.
\item
  \textbf{Use multiple temporal split points} (e.g., 60/40, 70/30,
  80/20) to assess sensitivity to the chosen split, and report mean ±
  std across splits.
\item
  \textbf{Always include Z-score and Isolation Forest as baselines.}
  These zero-parameter or minimal-parameter methods provide a lower
  bound that any proposed method should clear under honest evaluation.
\item
  \textbf{Separate detection from interpretation claims.} When a
  method's primary contribution is interpretability (e.g., root cause
  analysis, explanation generation) rather than raw detection accuracy,
  this should be explicitly stated and evaluated on
  interpretation-specific metrics.
\end{enumerate}

\subsubsection{5.3 Relationship to Existing
Approaches}\label{relationship-to-existing-approaches}

TSP is related to but distinct from existing time series validation
strategies:

\begin{itemize}
\tightlist
\item
  \textbf{Purged cross-validation} {[}36{]}, used in finance, removes
  training samples that overlap temporally with test samples. TSP
  eliminates the overlap at the source by separating the raw series
  before windowing.
\item
  \textbf{Blocked time series cross-validation} {[}37{]} splits time
  series into contiguous blocks but windows within blocks. TSP operates
  at the preprocessing stage, before any cross-validation strategy is
  applied.
\item
  \textbf{Walk-forward validation} {[}22{]} evaluates models
  sequentially on expanding windows. TSP provides a simpler, fixed-split
  alternative suitable for benchmarking where sequential retraining
  would be computationally prohibitive.
\end{itemize}

TSP's simplicity is intentional: it requires no new metrics, no new
benchmarks, and no additional infrastructure---only reordering two lines
of preprocessing code. Its adoption would immediately improve the
reliability of TSAD evaluation.

\subsubsection{5.4 Limitations of TSP}\label{limitations-of-tsp}

TSP makes an implicit stationarity assumption: that the data
distribution in the training period generalizes to the test period.
Under significant concept drift {[}38,39{]}, temporal-split AUC may
underestimate real-world performance if the deployed model is regularly
retrained. Conversely, if training data contains operational regimes
absent from the test period, temporal-split AUC may \emph{overestimate}
performance.

For datasets with very short duration or limited anomaly coverage in the
test period, temporal split may produce test sets with too few anomalies
for reliable metric computation. In such cases, we recommend reporting
both temporal-split results and explicitly noting the limitation, rather
than defaulting to random split.

TSP is a minimum standard, not a complete solution. It should be
combined with robust metrics (VUS-PR, not PA-F1), simple baselines, and
statistical significance testing.

\subsection{6. Discussion}\label{discussion}

\subsubsection{6.1 Reinterpreting the TSAD
Literature}\label{reinterpreting-the-tsad-literature}

If the mean leakage gap of ΔAUC = + observed in our experiments
generalizes, then many published TSAD results reporting AUC
\textgreater{} 0.90 on SWaT and similar benchmarks may correspond to
honest temporal-split AUCs in the 0.55--0.75 range. This does not mean
the published methods are worthless---many contribute genuinely novel
architectures and training strategies---but it does mean that \emph{the
field's perception of progress has been systematically inflated} by an
evaluation artifact.

This situation has a close historical parallel in finance. Before the
1990s, many quantitative trading strategies reported annual returns of
50\%+ in backtests---until the community recognized that look-ahead bias
(using information not available at the time of the trade) and
survivorship bias inflated results. The adoption of walk-forward
(temporal) validation and purged cross-validation corrected these
estimates and revealed that many ``profitable'' strategies were
artifacts of evaluation methodology {[}23,24{]}. TSAD appears to be at a
similar inflection point.

\subsubsection{6.2 Implications for the
Field}\label{implications-for-the-field}

\textbf{For researchers.} The most impactful contribution a TSAD paper
can make in 2026 may not be a new architecture, but a more honest
evaluation of existing architectures. We encourage the community to
revisit published methods under temporal-split evaluation and report
corrected performance estimates. The mTSBench {[}20{]} and TAB {[}19{]}
frameworks provide ready infrastructure for such re-evaluation.

\textbf{For benchmark maintainers.} We recommend that benchmark
platforms (TSB-AD, TAB, mTSBench, OpenTS-Bench) publish both
random-split and temporal-split baselines, and provide pre-defined
temporal split indices to ensure reproducible comparisons. Leaderboards
should display both numbers or default to temporal-split results.

\textbf{For conference reviewers and program chairs.} We recommend
requiring temporal-split evaluation and simple baseline comparisons
(Z-score, Isolation Forest) as a minimum standard for TSAD submissions,
analogous to the ablation study requirement that is already standard
practice.

\textbf{For industry practitioners.} The heterogeneity of leakage across
datasets (ΔAUC ranging from -0.254 to 0.104 in our study) means that
deployment decisions should be guided by dataset-specific temporal-split
evaluation, not by published random-split benchmark numbers. A method
that appears SOTA on a public leaderboard may underperform a simple
baseline on your specific data under honest evaluation.

\subsubsection{6.3 Limitations}\label{limitations}

Our study has several limitations. First, our five datasets all come
from cyber-physical systems and server/spacecraft monitoring domains.
Temporal leakage characteristics in other domains---finance, healthcare,
climate science, astronomy---may differ, and extending TSP evaluation to
these domains is important future work.

Second, our deep learning baselines (LSTM-AE, Transformer-AE) were
trained with fixed hyperparameters (15 epochs, learning rate 1e-3,
hidden size 64) for consistency and reproducibility. Hyperparameter
tuning per dataset and per method would likely improve absolute
performance, but the \emph{relative} comparison between random and
temporal split is unaffected.

Third, we did not evaluate the most recent foundation models (TimesFM,
MOIRAI, Chronos) or LLM-based methods due to computational constraints
(RTX 5060 8GB). Evaluating the temporal leakage characteristics of these
large-scale models---which may be more robust to temporal leakage due to
pre-training on diverse corpora---is an important direction for future
work.

Fourth, our overlap ablation is limited to three datasets and one method
(LSTM-AE), and the window length was fixed at L=256. Varying window
length independently of stride would provide a more complete picture of
the leakage mechanism.

\subsubsection{6.4 Beyond Detection: Future
Directions}\label{beyond-detection-future-directions}

Our finding that simple statistical methods match complex architectures
on current benchmarks does not mean TSAD research is complete. Rather,
it suggests that the field should shift focus from detection accuracy to
complementary capabilities:

\begin{itemize}
\tightlist
\item
  \textbf{Root cause analysis and anomaly localization}: identifying
  \emph{which} sensor(s) caused the anomaly and \emph{how} the fault
  propagated through the system {[}40,41{]}.
\item
  \textbf{Explanation generation}: producing natural-language
  explanations that operators can act on {[}42{]}.
\item
  \textbf{Open-set and zero-shot anomaly detection}: detecting anomaly
  types not seen during training, which requires moving beyond the
  closed-set assumption of current benchmarks {[}43,44{]}.
\item
  \textbf{Streaming and online evaluation}: deploying and evaluating
  detectors in continuous operation with human feedback {[}26{]}.
\end{itemize}

These capabilities require new benchmarks, new evaluation protocols, and
new metrics---all built on the foundation of honest temporal-split
evaluation that TSP provides.

\subsection{7. Conclusion}\label{conclusion}

We have demonstrated that temporal data leakage from random train/test
splits on overlapping sliding windows is a systematic, quantifiable
source of evaluation inflation in time series anomaly detection. Across
five datasets spanning cyber-physical systems, spacecraft telemetry, and
server monitoring, random splits inflate AUC-ROC by -0.009 on average.
The leakage magnitude is directly proportional to window overlap ratio,
confirming the causal mechanism. Under honest temporal-split evaluation,
simple statistical baselines match or exceed deep learning methods,
vision-language model features contribute zero additional detection
signal, and---critically---method rankings change, meaning that
random-split evaluation can lead to suboptimal model selection.

We have proposed the Temporal-Split Protocol (TSP)---a one-line
preprocessing change that splits time series before generating sliding
windows---as a minimum community evaluation standard. TSP requires no
new metrics, benchmarks, or infrastructure; it corrects a systematic
oversight in TSAD methodology; and its adoption would immediately
improve the reliability and comparability of published TSAD results. We
encourage the community to adopt TSP, re-evaluate published methods
under honest temporal splits, and direct future research toward the
capabilities that matter beyond detection accuracy: root cause analysis,
explanation generation, and robust deployment in real-world settings.

\subsection{References}\label{references}

{[}1{]} P. Malhotra et al., ``LSTM-based encoder-decoder for
multi-sensor anomaly detection,'' \emph{ICML 2016 Anomaly Detection
Workshop}, arXiv:1607.00148.

{[}2{]} J. Xu et al., ``Anomaly Transformer: Time Series Anomaly
Detection with Association Discrepancy,'' \emph{ICLR 2022},
arXiv:2110.02642.

{[}3{]} Y. Su et al., ``Robust Anomaly Detection for Multivariate Time
Series through Stochastic Recurrent Neural Network,'' \emph{KDD 2019},
pp.~2828--2837.

{[}4{]} H. Zhao et al., ``Multivariate Time-series Anomaly Detection via
Graph Attention Network,'' \emph{ICDM 2020}, pp.~841--850.

{[}5{]} Z. Darban et al., ``Deep Learning for Time Series Anomaly
Detection: A Survey,'' \emph{ACM Computing Surveys}, vol.~57, pp.~1--42,
2022.

{[}6{]} P. Malhotra et al., ``LSTM-based encoder-decoder for
multi-sensor anomaly detection,'' \emph{ICML 2016 Workshop},
arXiv:1607.00148.

{[}7{]} J. Xu et al., ``Anomaly Transformer,'' \emph{ICLR 2022}.

{[}8{]} S. Park et al., ``Delving into LLMs for Effective Time Series
Anomaly Detection,'' \emph{NeurIPS 2025}.

{[}9{]} Z. He et al., ``VLM4TS: Vision-Language Models for Time Series
Anomaly Detection,'' 2025.

{[}10{]} S. Kim et al., ``Towards a Rigorous Evaluation of Time-Series
Anomaly Detection,'' \emph{AAAI 2022}, vol.~36(7), pp.~7194--7201.

{[}11{]} M. Sehili and Z. Zhang, ``Multivariate Time Series Anomaly
Detection: Fancy Algorithms and Flawed Evaluation Methodology,''
\emph{TPCTC 2023}, LNCS vol.~14247, arXiv:2308.13068.

{[}12{]} M.S. Sarfraz et al., ``Position: Quo Vadis, Unsupervised Time
Series Anomaly Detection?'' \emph{ICML 2024}, arXiv:2405.02678.

{[}13{]} R. Wu and E. Keogh, ``Current Time Series Anomaly Detection
Benchmarks are Flawed and are Creating the Illusion of Progress,''
\emph{IEEE TKDE}, 2021.

{[}14{]} Lyu (2025). ``Did We Actually Fix It? An Independent
Adversarial Stress-Test of Post-Point-Adjustment Evaluation Metrics for
Time-Series Anomaly Detection.'' arXiv.

{[}15{]} A. Yugay et al., ``Strong Linear Baselines Strike Back,''
\emph{arXiv:2602.00672}, 2025.

{[}16{]} N. Xu, ``Trustworthiness and Timeseries Anomaly Detection
Methods,'' NAISS 2025/22-1491, KTH Royal Institute of Technology.

{[}17{]} ``Towards automated self-supervised learning for truly
unsupervised graph anomaly detection,'' \emph{Data Mining \& Knowledge
Discovery}, 2025.

{[}18{]} Liu \& Paparrizos, ``The Elephant in the Room: Towards A
Reliable Time-Series Anomaly Detection Benchmark,'' \emph{NeurIPS 2024}.

{[}19{]} Qiu et al., ``TAB: Unified Benchmarking of Time Series Anomaly
Detection Methods,'' \emph{PVLDB 2025}.

{[}20{]} ``mTSBench: Benchmarking Multivariate Time Series Anomaly
Detection and Model Selection at Scale,'' \emph{TMLR 2026}.

{[}21{]} ``TSB-AutoAD: Towards Automated Solutions for Time-Series
Anomaly Detection,'' \emph{SIGKDD 2025}.

{[}22{]} S.C. Hespeler et al., ``Temporal Cross-Validation Impacts
Multivariate Time Series Subsequence Anomaly Detection Evaluation,''
\emph{arXiv:2506.12183}, 2025.

{[}23{]} M. De Prado, \emph{Advances in Financial Machine Learning},
Wiley, 2018.

{[}24{]} D.H. Bailey et al., ``Pseudomathematics and Financial
Charlatanism,'' \emph{Notices of the AMS}, 61(5), 2014.

{[}25{]} Anonymous, ``Temporal Data Leakage from Pre-Split
Augmentation,'' \emph{arXiv:2512.06932}, 2025.

{[}26{]} ``Open Challenges in Time Series Anomaly Detection: An Industry
Perspective,'' \emph{arXiv}, 2025.

{[}27{]} ``NoBOOM: Real-World Chemical Process Anomaly Detection,''
\emph{NeurIPS 2025}.

{[}28{]} ``Beyond Academic Benchmarks: Critical Analysis and Best
Practices for Visual Industrial Anomaly Detection,'' \emph{CVPR 2025
Workshop}.

{[}29{]} J. Goh et al., ``A Dataset to Support Research in the Design of
Secure Water Treatment Systems,'' \emph{CRITIS 2016}.

{[}30{]} J.J. Downs and E.F. Vogel, ``A Plant-Wide Industrial Process
Control Problem,'' \emph{Computers \& Chemical Engineering}, 17(3),
1993.

{[}31{]} K. Hundman et al., ``Detecting Spacecraft Anomalies Using LSTMs
and Nonparametric Dynamic Thresholding,'' \emph{KDD 2018}.

{[}32{]} Y. Su et al., ``Robust Anomaly Detection for Multivariate Time
Series through Stochastic RNN,'' \emph{KDD 2019}.

{[}33{]} F.T. Liu et al., ``Isolation Forest,'' \emph{ICDM 2008}.

{[}34{]} B. Schölkopf et al., ``Estimating the Support of a
High-Dimensional Distribution,'' \emph{Neural Computation}, 2001.

{[}35{]} ViCSynAD project repository, 2026.

{[}36{]} M. De Prado, \emph{Advances in Financial Machine Learning},
Wiley, 2018. Chapter 7: Cross-Validation in Finance.

{[}37{]} C. Bergmeir and J.M. Benítez, ``On the use of cross-validation
for time series predictor evaluation,'' \emph{Information Sciences},
191, 2012.

{[}38{]} CALM: Concept drift adaptation via LLM-as-Judge for TSAD, 2025.

{[}39{]} DriftMind: Unsupervised concept drift detection, CPU-only,
2025.

{[}40{]} X. Han et al., ``Root Cause Analysis of Anomalies in
Multivariate Time Series through Granger Causal Discovery,'' \emph{ICLR
2025}.

{[}41{]} MultiverseAD: Spatial-temporal synchronous attention + causal
knowledge graph for TSAD, \emph{Neural Networks}, 2025.

{[}42{]} AXIS: Explainable Time Series Anomaly Detection with LLMs,
2025.

{[}43{]} DADA (ICLR 2025): Adaptive bottlenecks + dual adversarial
decoders for zero-shot TSAD.

{[}44{]} IBM TSPulse (ICLR 2026): 1M params, GPU-free, rank \#1 on
TSB-AD benchmark.

\printbibliography

\end{document}
